% Faithful translation baseline; Chinese source PDF pp. 53--61.
% Source retains Gaussian electromagnetic conventions (explicit c factors).
\section{Gauge Transformations}
\subsection*{1) Quantum coherence caused by a constant potential difference}
In classical mechanics, $\vect F=-\vect\nabla V(\vect x)$, so adding a spatially independent constant to the potential produces no effect.

In quantum mechanics, suppose a particle moves in a potential $V_1(\vect x)$. The evolution from the initial state $\ket{\alpha,t_0}$ is
\[\ket{\alpha,t}_1=e^{-\ii\op H(t-t_0)/\hbar}\ket{\alpha,t_0}.\]
If instead the particle moves in $V_2(\vect x,t)=V_1(\vect x)+\Delta V(t)$, the same initial state evolves as
\[\begin{split}
\ket{\alpha,t}_2
&=e^{-\ii\op H(t-t_0)/\hbar}
\exp\left[-\frac\ii\hbar\int_{t_0}^{t}\Delta V(t')\dd t'\right]\ket{\alpha,t_0}\\
&=\exp\left[-\frac\ii\hbar\int_{t_0}^{t}\Delta V(t')\dd t'\right]\ket{\alpha,t}_1.
\end{split}\]
The two states differ by a phase factor $e^{\ii(\varphi_2-\varphi_1)}$:
\[\varphi_2-\varphi_1=-\frac1\hbar\int_{t_0}^{t}\Delta V(t')\dd t'
\quad\xrightarrow{\Delta V\text{ constant}}\quad
-\frac{\Delta V}{\hbar}(t-t_0).\]
This may seem not to lead to any measurable result. However, interference between the two wavefunctions has a term $\cos(\varphi_2-\varphi_1)$ that depends on the phase difference and is observable.

\subsection*{2) Quantum coherence caused by a gravitational potential difference}
Consider the motion of thermal neutrons (uncharged) in a gravitational field.

At point A, the neutron beam is split into two: one travels through A--C--D to D, the other through A--B--D to D. If the frame ABCD is parallel to the ground, the two beams experience the same gravitational potential, and there is no relative gravitational phase at D.

\begin{center}\begin{tikzpicture}[>=Latex,scale=.95]
\draw[thick](0,0)rectangle(4.6,2.2);
\fill(0,0)circle(2pt)node[below left]{A};
\fill(0,2.2)circle(2pt)node[above left]{B};
\fill(4.6,0)circle(2pt)node[below right]{C};
\fill(4.6,2.2)circle(2pt)node[above right]{D};
\draw[<->](-.38,0)--(-.38,2.2)node[midway,left]{$l_2$};
\draw[<->](0,-.4)--(4.6,-.4)node[midway,below]{$l_1$};
\node[above]at(4.0,2.2){\scriptsize interference region};
\end{tikzpicture}\end{center}
If the frame is rotated upward by an angle $\delta$ around AC, the particle in the BD segment has a constant potential difference relative to AC:
\[\Delta V=mgl_2\sin\delta.\]
Although the AB and CD segments also have a gravitational potential difference from AC that depends on spatial position, their contributions are the same along both paths and do not affect the relative phase. Let the neutron's travel time from B to D be $T$. The resulting phase difference at D is
\[\varphi_{ABD}-\varphi_{ACD}=-\frac{m}{\hbar}gl_2T\sin\delta.\]
Since $T=l_1/v$, where $v$ is the neutron (wavepacket) velocity, and $mv=p=\hbar/\lambda$ with $\lambda$ the de Broglie wavelength in the source's convention, we have $T=l_1m\lambda/\hbar$ and
\[\boxed{\varphi_{ABD}-\varphi_{ACD}
=-\left(\frac{m}{\hbar}\right)^2gl_1l_2\lambda\sin\delta.}\]
As the angle $\delta$ changes, the interference term $\cos(\varphi_{ABD}-\varphi_{ACD})$ oscillates. Such oscillations have indeed been observed experimentally (see Fig.~2.6 in the source).

This discussion shows that a spatially uniform potential difference may leave the force unchanged while producing a quantum-interference effect.

\subsection*{3) Gauge transformations in electromagnetism}
Scalar and vector potentials:
\[\varphi(\vect x,t),\qquad\vect A(\vect x,t).\]
Fields:
\[\vect E=-\vect\nabla\varphi-\frac1c\frac{\partial\vect A}{\partial t},\qquad
\vect B=\vect\nabla\times\vect A.\]
From the Hamiltonian of classical electrodynamics, the quantum Hamiltonian of a charged particle in an electromagnetic field is
\[\boxed{\op H=\frac1{2m}\left(\vect{\op p}-\frac ec\vect A\right)^2+e\varphi,\quad
\varphi=\varphi(\vect{\op x},t),\ \vect A=\vect A(\vect{\op x},t).}\]
Using the fundamental commutation relations
\[[\op x_i,\op x_j]=[\op p_i,\op p_j]=0,\qquad
[\op x_i,\op p_j]=\ii\hbar\delta_{ij},\]
we have
\[\left[\op p_i-\frac ec A_i,\op p_j-\frac ec A_j\right]
=\frac{\ii\hbar e}{c}\epsilon_{ijk}B_k.\]
The Heisenberg equation of motion is
\[\begin{split}
\frac{\dd\op x_i}{\dd t}
&=\frac1{\ii\hbar}[\op x_i,\op H]\\
&=\frac1{2m\ii\hbar}
\left[\op x_i,\left(\op p_j-\frac ecA_j\right)
\left(\op p_j-\frac ecA_j\right)\right]\\
&=\frac1m\left(\op p_i-\frac ecA_i\right).
\end{split}\]
Consequently,
\[\begin{split}
m\frac{\dd^2\op x_i}{\dd t^2}
&=\frac{\dd}{\dd t}\left(\op p_i-\frac ecA_i\right)\\
&=\frac1{\ii\hbar}\left[\op p_i-\frac ecA_i,\op H\right]
-\frac ec\frac{\partial A_i}{\partial t}\\
&=\frac{e}{2mc}\epsilon_{ijk}
\left(\op p_j-\frac ecA_j\right)B_k
-\frac{e}{2mc}\epsilon_{ijk}B_j\left(\op p_k-\frac ecA_k\right)
+eE_i.
\end{split}\]
In vector notation,
\[\boxed{m\frac{\dd^2\vect{\op x}}{\dd t^2}
=e\left[\vect E+\frac1{2c}
\left(\frac{\dd\vect{\op x}}{\dd t}\times\vect B
-\vect B\times\frac{\dd\vect{\op x}}{\dd t}\right)\right].}\]
This is the Ehrenfest theorem in the Heisenberg picture. Note that the velocity operator $\dd\vect{\op x}/\dd t$ and $\vect B=\vect\nabla\times\vect A$ generally do not commute; operator ordering matters in this expression.

In the coordinate representation of the Schr\"odinger picture,
\[\ii\hbar\frac{\partial\psi(\vect x,t)}{\partial t}
=\left[\frac1{2m}\left(-\ii\hbar\vect\nabla-\frac ec\vect A(\vect x,t)\right)^2
+e\varphi(\vect x,t)\right]\psi(\vect x,t).\]
The continuity equation is
\[\frac{\partial\rho(\vect x,t)}{\partial t}
+\vect\nabla\cdot\vect j(\vect x,t)=0,\]
where
\[\rho(\vect x,t)=|\psi(\vect x,t)|^2,\]
\[\vect j(\vect x,t)=\frac\hbar m\operatorname{Im}
\left[\psi^*(\vect x,t)
\left(\vect\nabla-\frac{\ii e}{\hbar c}\vect A(\vect x,t)\right)
\psi(\vect x,t)\right].\]
Now transform the potentials as
\[\varphi(\vect x,t)\longrightarrow\varphi(\vect x,t),\qquad
\vect A(\vect x,t)\longrightarrow\vect A(\vect x,t)+\vect\nabla\Lambda(\vect x),\]
where $\Lambda(\vect x)$ is an arbitrary scalar function. It is clear that the fields $\vect E$ and $\vect B$ remain unchanged. A transformation of potentials leaving the field strengths unchanged is called a \emph{gauge transformation}; distinct functions $\Lambda(\vect x)$ specify different gauges.

Let the state corresponding to the potentials $\varphi,\vect A$ be $\ket{\psi,t}_1$, and that corresponding to $\varphi,\vect A+\vect\nabla\Lambda$ be $\ket{\psi,t}_2$. They satisfy respectively
\[\ii\hbar\frac{\partial}{\partial t}\ket{\psi,t}_1
=\left[\frac{(\vect{\op p}-e\vect A/c)^2}{2m}+e\varphi\right]\ket{\psi,t}_1,\]
\[\ii\hbar\frac{\partial}{\partial t}\ket{\psi,t}_2
=\left[\frac{[\vect{\op p}-e(\vect A+\vect\nabla\Lambda)/c]^2}{2m}
+e\varphi\right]\ket{\psi,t}_2.\]
Take
\[\ket{\psi,t}_2=\exp\left(\frac{\ii e}{\hbar c}\Lambda(\vect{\op x})\right)\ket{\psi,t}_1.\]
Because
\[e^{-\ii e\Lambda/(\hbar c)}
\left[\vect{\op p}-\frac ec(\vect A+\vect\nabla\Lambda)\right]
e^{\ii e\Lambda/(\hbar c)}
=\vect{\op p}-\frac ec\vect A,\]
it is easy to show that the two Schr\"odinger equations are consistent. Thus, under a gauge transformation the states differ by a phase factor:
\[\boxed{\ket{\psi,t}\longrightarrow
\exp\left(\frac{\ii e}{\hbar c}\Lambda(\vect{\op x})\right)\ket{\psi,t}.}\]
Does a gauge transformation have observable quantum effects?

\subsection*{4) Aharonov--Bohm (AB) effect}
Consider a long hollow cylindrical shell with inner radius $\rho_a$ and outer radius $\rho_b$. The particle moves in the shell $\rho_a<\rho<\rho_b$; it cannot enter $\rho<\rho_a$ or $\rho>\rho_b$. The wavefunction vanishes at $\rho=\rho_a$ and $\rho=\rho_b$, so this is a bound-state problem.

Now place a long solenoid in the central hollow region, with its magnetic field confined inside the solenoid, ensuring that the field outside vanishes.
\begin{center}
\begin{tikzpicture}[>=Latex,scale=.75]
\begin{scope}
\draw (0,1.4) ellipse (1.55 and .48);
\draw (0,-1.4) arc[start angle=180,end angle=360,x radius=1.55,y radius=.48];
\draw (-1.55,-1.4)--(-1.55,1.4);
\draw (1.55,-1.4)--(1.55,1.4);
\draw[dashed] (0,1.4) ellipse (.65 and .20);
\draw[dashed] (-.65,-1.4)--(-.65,1.4);
\draw[dashed] (.65,-1.4)--(.65,1.4);
\draw[->] (0,1.4)--(1.0,1.78) node[right]{$\rho_b$};
\draw[->] (0,1.4)--(.55,1.55) node[above]{$\rho_a$};
\draw[<->] (1.9,-1.4)--(1.9,1.4) node[midway,right]{$L$};
\end{scope}
\begin{scope}[xshift=6cm]
\draw (0,1.4) ellipse (1.55 and .48);
\draw (0,-1.4) arc[start angle=180,end angle=360,x radius=1.55,y radius=.48];
\draw (-1.55,-1.4)--(-1.55,1.4);
\draw (1.55,-1.4)--(1.55,1.4);
\draw[dashed] (0,1.4) ellipse (.65 and .20);
\draw[dashed] (0,-1.4) ellipse (.65 and .20);
\foreach \x in {-.48,-.18,.18,.48}{
\draw[->] (\x,-1.85)--(\x,2.0);}
\end{scope}
\end{tikzpicture}
\end{center}
Assume
\[\vect B=\begin{cases}0,&\rho>\rho_a,\\B\vect e_z,&\rho<\rho_a.\end{cases}\]
In cylindrical coordinates,
\[\vect B=\vect\nabla\times\vect A
=\frac1\rho\begin{vmatrix}
\vect e_\rho&\rho\vect e_\phi&\vect e_z\\
\partial/\partial\rho&\partial/\partial\phi&\partial/\partial z\\
A_\rho&\rho A_\phi&A_z
\end{vmatrix}.\]
One finds
\[\boxed{\vect A=\begin{cases}
\dfrac{B\rho_a^2}{2\rho}\vect e_\phi,&\rho>\rho_a,\\[6pt]
\dfrac{B\rho}{2}\vect e_\phi,&\rho<\rho_a.
\end{cases}}\]
Introducing the field into the region $\rho<\rho_a$ changes the vector potential in the particle's accessible region $\rho_a<\rho<\rho_b$ from $\vect A=0$ to $\vect A\ne0$, while the magnetic field there remains $\vect B=0$.

The stationary Schr\"odinger equation is
\[\frac1{2m}\left(\vect{\op p}-\frac ec\vect A\right)^2\psi
=-\frac{\hbar^2}{2m}\left[
\vect e_\rho\frac{\partial}{\partial\rho}
+\frac{\vect e_\phi}{\rho}\left(\frac{\partial}{\partial\phi}
-\frac{\ii e B\rho_a^2}{2\hbar c}\right)
+\vect e_z\frac{\partial}{\partial z}\right]^2\psi
=E\psi,\quad\rho_a<\rho<\rho_b.\]
Clearly, the magnetic-field term changes the particle's bound-state energies $E$. Crucially, the magnetic field $\vect B$ is absent from the particle's accessible region, while $\vect A$ is present there. In classical mechanics the motion of a charged particle is determined by $\vect E$ and $\vect B$, whereas in quantum mechanics a potential without a field strength can still influence the properties of a system. This is a purely quantum effect.

The above is the influence of the AB effect on a bound-state problem.

Now consider the AB effect using the path integral. A charged particle starts at A and passes through two slits in a screen to reach B, where interference occurs. A long solenoid lies between the two slits behind the screen. The magnetic field is inside the solenoid and vanishes outside. For the classical system, the field inside the solenoid cannot affect the interference at B.

\begin{center}\begin{tikzpicture}[>=Latex,scale=.85]
\fill (0,0)circle(2pt)node[left]{A};\fill (6,0)circle(2pt)node[right]{B};
\draw[thick](2,-1.3)--(2,-.8);\draw[thick](2,-.56)--(2,.56);\draw[thick](2,.8)--(2,1.3);
\draw[thick](4,-1.3)--(4,1.3);
\draw[thick] (3,0)ellipse(.65 and .40);\node at(3,0){$\vect B\ne0$};
\draw[thick,->](0,0)to[out=33,in=160] (2,.70)to[out=-8,in=160](6,0);
\draw[thick,->](0,0)to[out=-33,in=200](2,-.70)to[out=8,in=200](6,0);
\node[above]at(3.2,1.25){\scriptsize impenetrable cylinder};
\end{tikzpicture}\end{center}
For a quantum system, although the field vanishes outside the solenoid, $\vect A\ne0$. The action in the path integral is written in the source as
\[\begin{split}
S[\vect x(t)]
&=\int_{t_A}^{t_B}\dd t\,\frac1{2m}\left(\vect p-\frac ec\vect A\right)^2\\
&\simeq\int_{t_A}^{t_B}\dd t\left[
\frac m2\left(\frac{\dd\vect x}{\dd t}\right)^2
-\frac ec\frac{\dd\vect x}{\dd t}\cdot\vect A\right]\\
&=S_0[\vect x(t)]-\frac ec\int_{t_A}^{t_B}\dd t\,
\frac{\dd\vect x}{\dd t}\cdot\vect A
=S_0[\vect x(t)]-\frac ec\int_{\vect x_A}^{\vect x_B}\vect A\cdot\dd\vect x.
\end{split}\]
The source neglects the term quadratic in $\vect A$. For any path above the solenoid and any path below it, the two contributions to the amplitude $K(\vect x_B,t_B;\vect x_A,t_A)$ have the form
\[\begin{split}
&\exp\left[-\frac{\ii e}{\hbar c}\int_{\vect x_A,\,\mathrm{above}}^{\vect x_B}\vect A\cdot\dd\vect x\right]
 e^{\ii S_0[\vect x_{\mathrm{above}}(t)]/\hbar}\\
&\quad+\exp\left[-\frac{\ii e}{\hbar c}\int_{\vect x_A,\,\mathrm{below}}^{\vect x_B}\vect A\cdot\dd\vect x\right]
 e^{\ii S_0[\vect x_{\mathrm{below}}(t)]/\hbar}.
\end{split}\]
There is a relative phase that affects the interference fringes:
\[\begin{split}
\frac e{\hbar c}\left[
\int_{\mathrm{above}}\vect A\cdot\dd\vect x
-\int_{\mathrm{below}}\vect A\cdot\dd\vect x\right]
&=\frac e{\hbar c}\oint\vect A\cdot\dd\vect x\\
&=\frac e{\hbar c}\int(\vect\nabla\times\vect A)\cdot\dd\vect S\\
&=\frac e{\hbar c}\int\vect B\cdot\dd\vect S
=\frac e{\hbar c}B\pi R^2.
\end{split}\]
This is independent of the precise paths; Stokes' theorem has been used.

The AB effect shows that in quantum mechanics, even when the field strength vanishes, a nonzero potential can influence the interference of particles.

\subsection*{5) Magnetic monopoles}
In classical electromagnetic theory, static electric and magnetic fields are asymmetric. The static electric field has sources but no curl:
\[\vect\nabla\cdot\vect E=4\pi\rho_e,\qquad
\vect\nabla\times\vect E=0.\]
Electric field-line patterns (single-charge configuration):
\begin{center}\begin{tikzpicture}[>=Latex,scale=.62]
\fill (0,0)circle(10pt)node[white]{\small $+$};
\foreach \a in {0,45,...,315}{\draw[->] (\a:0.25)--(\a:1.5);}
\begin{scope}[xshift=6cm]
\fill (0,0)circle(10pt)node[white]{\small $-$};
\foreach \a in {0,45,...,315}{\draw[->] (\a:1.5)--(\a:.25);}
\end{scope}
\end{tikzpicture}\end{center}
The static magnetic field has no sources and does have curl:
\[\vect\nabla\cdot\vect B=0,\qquad
\vect\nabla\times\vect B=\frac{4\pi}{c}\vect J.\]
The magnetic field-line pattern has a dipole structure:
\begin{center}\begin{tikzpicture}[>=Latex,scale=.65]
\draw[thick,fill=gray!12](-1,-.25)rectangle(1,.25);
\node at(-.72,0){N};\node at(.72,0){S};
\foreach \h in {.45,.75,1.12,1.55}{
\draw[->,smooth]plot coordinates{(-.8,.25)(-1.1,\h)(0,\h+0.5)(1.1,\h)(.8,.25)};
\draw[->,smooth]plot coordinates{(-.8,-.25)(-1.1,-\h)(0,-\h-.5)(1.1,-\h)(.8,-.25)};}
\end{tikzpicture}\end{center}
Question: Can there be an isolated magnetic charge, that is, a magnetic monopole?

\paragraph{Dirac string.}
If there is a magnetic charge $e_m$ at the origin, the static field it produces satisfies
\[\vect\nabla\cdot\vect B=4\pi e_m\delta^{(3)}(\vect r).\]
Its solution is
\[\boxed{\vect B(\vect r)=\frac{e_m}{r^2}\vect e_r.}\]
Can a vector potential $\vect A(\vect r)$ be found such that the original relation $\vect B=\vect\nabla\times\vect A$ still holds everywhere? The answer is no.

Proof: If $\vect B=\vect\nabla\times\vect A$, then necessarily
\[\vect\nabla\cdot\vect B=\vect\nabla\cdot(\vect\nabla\times\vect A)=0,\]
contradicting the monopole source equation.

Likewise, by Gauss's theorem, for a closed surface surrounding the monopole,
\[\oint\vect B\cdot\dd\vect S=4\pi e_m.\]
But if $\vect B=\vect\nabla\times\vect A$, then
\[\oint\vect B\cdot\dd\vect S
=\oint(\vect\nabla\times\vect A)\cdot\dd\vect S
=\int\vect\nabla\cdot(\vect\nabla\times\vect A)\dd V=0,\]
contradicting Gauss's theorem. Thus no single globally regular vector potential $\vect A(\vect r)$ can produce $\vect B(\vect r)=\vect\nabla\times\vect A$.

For example, in spherical coordinates take
\[\vect A=e_m\frac{1-\cos\theta}{r\sin\theta}\vect e_\varphi.\]
For $\theta<\pi-\varepsilon$ ($\varepsilon$ infinitesimal), using
\[\begin{split}
\vect\nabla\times\vect A
&=\frac1{r\sin\theta}\left[
\frac{\partial}{\partial\theta}(A_\varphi\sin\theta)
-\frac{\partial A_\theta}{\partial\varphi}\right]\vect e_r\\
&\quad+\frac1r\left[
\frac1{\sin\theta}\frac{\partial A_r}{\partial\varphi}
-\frac{\partial}{\partial r}(rA_\varphi)\right]\vect e_\theta\\
&\quad+\frac1r\left[
\frac{\partial}{\partial r}(rA_\theta)
-\frac{\partial A_r}{\partial\theta}\right]\vect e_\varphi
=\frac{e_m}{r^2}\vect e_r=\vect B.
\end{split}\]
However, $\vect A$ diverges at $\theta=\pi$. The negative $z$ axis therefore contains a singular string, called the Dirac string. Its singular contribution means that the distributional curl of this globally extended singular potential does not equal the monopole field alone.

\paragraph{Quantization of electric charge.}
Although a single vector potential satisfying $\vect B=\vect\nabla\times\vect A$ cannot cover all space, such potentials can be found in separate regions. For example, take
\[\boxed{\begin{aligned}
\vect A_{\mathrm I}&=e_m\frac{1-\cos\theta}{r\sin\theta}\vect e_\varphi,
\qquad\theta<\pi-\varepsilon,\\
\vect A_{\mathrm{II}}&=-e_m\frac{1+\cos\theta}{r\sin\theta}\vect e_\varphi,
\qquad\theta>\varepsilon.
\end{aligned}}\]
$\vect A_{\mathrm I}$ diverges at $\theta=\pi$, while $\vect A_{\mathrm{II}}$ diverges at $\theta=0$. With this choice of potentials,
\[\vect\nabla\times\vect A=\frac{e_m}{r^2}\vect e_r=\vect B\]
in each nonsingular region.

In their overlap, since both potentials produce the same magnetic field, they are related by a gauge transformation $\Lambda$:
\[\vect A_{\mathrm{II}}-\vect A_{\mathrm I}
=-\frac{2e_m}{r\sin\theta}\vect e_\varphi=\vect\nabla\Lambda.\]
In spherical coordinates,
\[\vect\nabla\Lambda=\frac{\partial\Lambda}{\partial r}\vect e_r
+\frac1r\frac{\partial\Lambda}{\partial\theta}\vect e_\theta
+\frac1{r\sin\theta}\frac{\partial\Lambda}{\partial\varphi}\vect e_\varphi.\]
Hence
\[\boxed{\Lambda(\vect r)=-2e_m\varphi.}\]
Consider a particle of electric charge $e$ moving in the monopole field. The wavefunctions $\psi_{\mathrm I}$ and $\psi_{\mathrm{II}}$ corresponding to $\vect A_{\mathrm I}$ and $\vect A_{\mathrm{II}}$ satisfy in the overlapping region
\[\psi_{\mathrm{II}}=e^{\ii e\Lambda/(\hbar c)}\psi_{\mathrm I}
=e^{-2\ii ee_m\varphi/(\hbar c)}\psi_{\mathrm I}.\]
By single-valuedness of the wavefunction, $\psi_{\mathrm{II}}$ must be single valued in the overlap. For example, at $\theta=\pi/2$,
\[\psi_{\mathrm{II}}(r,\theta=\pi/2,\varphi=0)
=\psi_{\mathrm{II}}(r,\theta=\pi/2,\varphi=2\pi).\]
Consequently,
\[\frac{2ee_m}{\hbar c}=\pm n,\qquad n=0,\pm1,\pm2,\ldots.\]
This is electric-charge quantization:
\[\boxed{e=\pm\frac{\hbar c}{2e_m}n.}\]
There is a smallest unit of electric charge.

Quantum mechanics does not require the existence of magnetic monopoles. If a magnetic monopole exists, however, electric charge must be quantized and there is a smallest unit of electric charge.